\section{Predictive Modeling Design}
\label{sec:modeling_design}

The predictors retained in Section~\ref{sec:feature_engineering} define three nested tabular information sets: market features; market and conventional ownership features; and market, ownership, and engineered network features. Graph neural networks instead receive the retained market and ownership variables as stock-node attributes, but not the engineered network features, because they learn directly from the manager--stock graph and its ownership edges.

The predictive analysis proceeds through a sequence of validation-based comparisons. First, the three nested feature sets are evaluated separately within Ridge and XGBoost under common initial settings. Comparisons across feature sets test whether conventional ownership and engineered network features provide incremental predictive information (Hypotheses~H1 and~H2), while comparisons between Ridge and XGBoost on identical inputs test whether nonlinear flexibility improves prediction (Hypothesis~H3). One feature specification is retained within each model class. When validation differences are negligible, the simpler information set is preferred.

Second, the retained XGBoost specification is refined using the controlled tuning procedure described below, while Ridge retains its fixed penalty. Third, GraphSAGE, GAT, and Temporal GraphSAGE are compared under a common initial configuration, after which only the strongest graph architecture is tuned. Comparing the tuned graph specification with the tabular finalists tests whether learning directly from the ownership structure provides incremental predictive value (Hypothesis~H4).

The final validation comparison considers the retained linear, nonlinear, and graph-based specifications. Model selection is completed before the test sample is examined and considers the primary validation metric together with the consistency of the complementary metrics and model complexity. Additional complexity is retained only when it produces a material and consistent predictive improvement. Once the decision is fixed, the finalist specifications are refitted on the combined training and validation samples and evaluated once on the held-out test sample. Test performance is used to assess whether the validation evidence generalizes, not to reopen model selection.

\subsection{Chronological Samples and Model Selection}
\label{sec:chronological_samples}

After manager-quarter filings are consolidated, securities are mapped to historical CRSP records, and observations are checked for usable predictors and outcomes, eligible stock-quarters are assigned to fixed chronological training, validation, and test samples. Table~\ref{tab:modeling_samples} reports their sizes and coverage. Target distributions across the chronological modeling samples are reported in Appendix~\ref{app:chronological_samples}.

\begin{table}[H]
    \centering
    \caption{Chronological modeling samples}
    \label{tab:modeling_samples}
    \begin{tabular}{lrrl}
        \toprule
        Sample & Observations & Quarters & Coverage \\
        \midrule
        Training   & 123,791 & 34 & 2013 Q2--2021 Q3 \\
        Validation & 29,716 & 7 & 2022 Q1--2023 Q3 \\
        Test       & 30,511 & 8 & 2024 Q1--2025 Q4 \\
        \bottomrule
    \end{tabular}
\end{table}


\noindent
Chronological validation is used instead of random assignment because the empirical objective is to predict later quarters from earlier information. Random assignment would mix observations from different periods and would not reproduce this forward-looking setting. Because the primary target uses the 63 trading days after each information cutoff, 2021~Q4 and 2023~Q4 are purged at the training--validation and validation--test boundaries, respectively. This ensures that the outcomes used in one sample are resolved before the following sample begins.

The training sample is used to estimate candidate models. The validation sample is used to compare specifications, tune hyperparameters, and determine the number of boosting rounds or training epochs. During this process, the preprocessing parameters described in Section~\ref{sec:feature_engineering} are estimated from the training sample and applied unchanged to validation observations. Validation MAE is the primary selection metric. When competing specifications achieve materially similar MAE, the complementary metrics defined in Section~\ref{sec:evaluation_metrics}, together with model complexity, are used to prefer the more stable and parsimonious specification.

Once model selection has been completed using validation data, the finalist specifications are refitted on the combined training and validation samples. Their feature sets and hyperparameters, including the validation-selected number of boosting rounds or training epochs where applicable, remain fixed. Each finalist is then evaluated once on the held-out test sample to assess whether the validation evidence generalizes. Test performance does not reopen model selection.

\subsection{No-Information Benchmark}
\label{sec:no_information_benchmark}

The no-information benchmark predicts the fitting-sample mean target for every observation. For fitting sample \(\mathcal{F}\), its prediction \(\hat{y}_i\) is

\begin{equation}
    \hat{y}^{\,\mathrm{naive}}_i
    =
    \bar{y}_{\mathcal{F}}
    =
    \frac{1}{N_{\mathcal{F}}}
    \sum_{j\in\mathcal{F}} y_j,
    \label{eq:naive_mean}
\end{equation}

\noindent
where \(N_{\mathcal{F}}\) is the number of fitting observations. The benchmark uses no stock-level predictors and provides a direct reference for whether a fitted model improves on a constant forecast. For validation, \(\mathcal{F}\) is the training sample; for the final test comparison, it is the combined training and validation sample.

\subsection{Ridge Regression}
\label{sec:ridge_regression}

The linear model is Ridge regression \parencite{hoerl1970ridge}. Its coefficients solve

\begin{equation}
    \left(\hat{\beta}_0,\hat{\boldsymbol{\beta}}\right)
    =
    \arg\min_{\beta_0,\boldsymbol{\beta}}
    \left\{
    \sum_{i=1}^{N}
    \left(
    y_i-\beta_0-\mathbf{x}_i^{\prime}\boldsymbol{\beta}
    \right)^2
    +
    \alpha\lVert\boldsymbol{\beta}\rVert_2^2
    \right\},
    \label{eq:ridge_objective}
\end{equation}

\noindent
where \(\mathbf{x}_i\) contains the standardized predictors, \(\boldsymbol{\beta}\) contains their coefficients, and the intercept \(\beta_0\) is not penalized. The \(L_2\) penalty stabilizes coefficient estimates when predictors remain related after feature selection, while standardization makes the penalty comparable across variables measured on different scales. For simplicity, \(\alpha\) is fixed at 1 across all Ridge specifications. Holding the penalty constant ensures that differences across feature sets reflect changes in the information supplied to the model rather than simultaneous changes in regularization.

\subsection{XGBoost}
\label{sec:nonlinear_models}

The nonlinear tabular model is an XGBoost gradient-boosted tree ensemble \parencite{chen2016xgboost}. It builds predictions sequentially, with each additional tree seeking to reduce the errors left by the preceding trees. This structure allows nonlinear relationships and interactions among predictors to be learned without specifying them in advance.

A common prespecified configuration is used for the initial feature-set comparisons. The retained XGBoost specification is then refined through a controlled validation procedure that considers changes in learning rate and tree complexity, while the number of boosting rounds is determined through validation-based early stopping. XGBoost is trained using a squared-error objective, and validation MAE is the primary tuning metric. When configurations perform similarly, the faster-converging and less complex alternative is preferred. The complete initial settings and tuning configurations are reported in Appendix~\ref{app:model_specifications}.

\subsection{Graph-Based Models}
\label{sec:graph_models}

The graph-based analysis uses the quarterly manager--stock networks defined in Section~\ref{sec:bipartite_graph}. Stock nodes receive the retained market and conventional ownership features, while manager nodes receive portfolio value, security count, largest position weight, and portfolio concentration. Ownership edges carry portfolio weights and reported-ownership shares. The GNNs do not receive the engineered network features used by the tabular network-augmented specification, since they learn directly from the graph and its ownership edges.

Three architectures are considered. Weighted GraphSAGE \parencite{hamilton2017inductive} aggregates neighboring nodes using observed ownership-edge weights, while GAT \parencite{velickovic2018graph} learns attention weights from node and edge information. Both process each quarterly graph independently. Temporal GraphSAGE additionally uses a gated recurrent unit to carry stock-node states across consecutive quarters.

Figure~\ref{fig:weighted_graphsage_update} illustrates the principal elements of a manager-to-stock update in weighted GraphSAGE. The representation is conceptual; the model-specific message-passing and temporal-update equations are provided in Appendix~\ref{app:gnn_architectures}.


\begin{figure}[!htbp]
    \centering
    \includegraphics[width=\textwidth]{08_05_gnn_node_update.pdf}
    \caption[Weighted GraphSAGE message passing]{Weighted GraphSAGE update for a stock node using normalized ownership weights $a_{mi}\equiv o_{mi,q}$, with $\sum_{m\in\mathcal{N}(i)}a_{mi}=1$, followed by a learned transformation and nonlinear activation.}
    \label{fig:weighted_graphsage_update}
\end{figure}

\noindent
The three architectures are initially estimated under a common configuration and compared using validation MAE. Only the strongest architecture is subsequently refined through controlled changes in model capacity, regularization, and learning rate, while the number of training epochs is determined through validation-based early stopping. The resulting graph specification then enters the final validation comparison with the retained linear and nonlinear tabular specifications. Overall model selection follows the performance and parsimony criteria described in Section~\ref{sec:chronological_samples}.

The graph-modeling sequence contains 51 labeled snapshots because the final ownership snapshot has no resolved 63-day outcome. Purged snapshots may update temporal states but contribute no training loss. Complete model settings and tuning configurations are reported in Appendix~\ref{app:model_specifications}.

\subsection{Evaluation Metrics}
\label{sec:evaluation_metrics}

Every model predicts the continuous downside-risk target defined in Section~\ref{sec:target_definition}. Predictions are bounded between 0 and 1 before evaluation, corresponding to the feasible range of the maximum-drawdown target. Let \(y_i\) denote the realized target and \(\hat{y}_i\) its prediction for observation \(i\) in an evaluation sample of size \(N\). The primary quantitative selection metric is the mean absolute error,

\begin{equation}
    \mathrm{MAE}
    =
    \frac{1}{N}
    \sum_{i=1}^{N}
    \left|y_i-\hat{y}_i\right|.
    \label{eq:mae}
\end{equation}

\noindent
MAE expresses prediction error in the same units as the target and is less dominated by the target's severe upper tail than a squared-error criterion.

Two complementary magnitude-based metrics are also reported. Root mean squared error gives greater weight to large errors,

\begin{equation}
    \mathrm{RMSE}
    =
    \sqrt{
        \frac{1}{N}
        \sum_{i=1}^{N}
        \left(y_i-\hat{y}_i\right)^2
    },
    \label{eq:rmse}
\end{equation}

\noindent
while the coefficient of determination compares squared prediction error with a constant prediction equal to the evaluation-sample mean,

\begin{equation}
    R^2
    =
    1-
    \frac{
        \sum_{i=1}^{N}
        \left(y_i-\hat{y}_i\right)^2
    }{
        \sum_{i=1}^{N}
        \left(y_i-\bar{y}\right)^2
    },
    \label{eq:r_squared}
\end{equation}

\noindent
where \(\bar{y}\) is the mean realized target in that sample. An \(R^2\) below zero therefore indicates that the model performs worse under squared error than this constant benchmark.

For the Ridge feature-set comparison, the Akaike information criterion (AIC) and Bayesian information criterion (BIC) are additionally reported as in-sample complexity diagnostics. Because Ridge shrinks rather than eliminates coefficients, model complexity is measured using the effective degrees of freedom,

\begin{equation}
    d_{\alpha}
    =
    1+
    \sum_k
    \frac{s_k^2}{s_k^2+\alpha},
    \label{eq:ridge_effective_df}
\end{equation}

\noindent
where \(s_k\) are the singular values of the standardized design matrix and the additional unit represents the intercept. Let \(n\) denote the number of training observations and \(\mathrm{RSS}\) the training residual sum of squares. The information criteria are calculated as

\begin{equation}
    \begin{aligned}
        \mathrm{AIC}
         & =
        n\left[
            \log(2\pi)+1+
            \log\left(\frac{\mathrm{RSS}}{n}\right)
            \right]
        +2d_{\alpha}, \\
        \mathrm{BIC}
         & =
        n\left[
            \log(2\pi)+1+
            \log\left(\frac{\mathrm{RSS}}{n}\right)
            \right]
        +d_{\alpha}\log(n).
    \end{aligned}
    \label{eq:ridge_information_criteria}
\end{equation}

\noindent
These statistics are reported only for the linear feature-set comparison and do not determine model selection.

Because the economic analysis uses predictions to rank stocks within each quarter, ranking performance is evaluated separately. For every evaluation quarter \(q\), the Spearman correlation \(\rho_q\) is computed between predicted and realized downside risk across the \(N_q\) eligible stocks. The reported statistic is the equally weighted mean across the \(Q\) evaluated quarters,

\begin{equation}
    \overline{\rho}
    =
    \frac{1}{Q}
    \sum_{q=1}^{Q}
    \rho_q.
    \label{eq:mean_quarterly_spearman}
\end{equation}

\noindent
A value near zero indicates little systematic within-quarter ranking ability, whereas a positive value indicates that stocks assigned greater predicted fragility tend to experience larger subsequent drawdowns. This metric directly evaluates the within-quarter ranking that underlies the Fragility Score defined in Section~\ref{sec:fragility_score_construction}.

If XGBoost is selected, SHAP (SHapley Additive exPlanations) decomposes each fitted prediction into additive feature contributions relative to the model's expected output, summarizing their global importance and observation-level direction \parencite{lundberg2017unified}.
